\section{Poisson likelihood under misspecification}
Write $n$ for the sample size in the asymptotic arguments;
the stated sample has $n=1000$.

\subsection*{(a) Correctly specified Poisson model}
The log likelihood is
\[
\ell_n(\lambda)=-n\lambda+
\left(\sum_{i=1}^ny_i\right)\log\lambda-\sum_{i=1}^n\log(y_i!).
\]
Setting its derivative to zero and using invariance of maximum
likelihood gives
\[
\boxed{\widehat\lambda_{\mathrm{ML}}=\bar y,\qquad
\widehat P_0=e^{-\bar y}.}
\]
For an interior value $\lambda>0$,
$\sqrt n(\bar y-\lambda)\xrightarrow{d}N(0,\lambda)$.
The delta method, with $h'(\lambda)=-e^{-\lambda}$, then gives
\[
\boxed{\sqrt n(\widehat P_0-e^{-\lambda})
\xrightarrow{d}N(0,\lambda e^{-2\lambda}).}
\]
For $n=1000$, the approximate variance of $\widehat P_0$ is
$\lambda e^{-2\lambda}/1000$, consistently estimated by
$\bar y e^{-2\bar y}/1000$. If all observations equal zero, the MLE
is the boundary value $\widehat\lambda=0$ when this is included
in the parameter space.

\subsection*{(b) Negative binomial data}
For a proper nondegenerate $NB(r,p)$ distribution, take $r\geq1$
and $0<p<1$. Its mean and variance are
\[
\mu=\frac{r(1-p)}p,\qquad v=\frac{r(1-p)}{p^2}.
\]
The law of large numbers and continuous mapping theorem imply
\[
\boxed{\widehat\lambda_{\mathrm{ML}}\xrightarrow{p}
\lambda^*=\frac{r(1-p)}p,\qquad
\widehat P_0\xrightarrow{p}
\exp\left[-\frac{r(1-p)}p\right].}
\]
The actual zero probability is $P_0=p^r$.
For $0<p<1$,
\[
\log p>-\frac{1-p}{p}
\quad\Longrightarrow\quad
\exp\left[-\frac{r(1-p)}p\right]<p^r.
\]
Thus the fitted Poisson model asymptotically understates the zero
probability. At $p=1$, both distributions degenerate at zero and both
zero probabilities equal one. The endpoint $p=0$ does not define a
proper finite-valued distribution with the pmf in the question.

\clearpage
\subsection*{(c) Kullback--Leibler minimizing parameter}
Collecting all terms that do not depend on $\lambda$ in a constant $K$,
\[
d(g_{r,p},f_\lambda)=K+\lambda-\mu\log\lambda.
\]
For $0<p<1$,
\[
\frac{\partial d}{\partial\lambda}=1-\frac\mu\lambda,\qquad
\frac{\partial^2d}{\partial\lambda^2}=\frac\mu{\lambda^2}>0.
\]
Hence the unique minimizing value is
\[
\boxed{\lambda^*=\mu=\frac{r(1-p)}p.}
\]
This equals the probability limit in (b). The misspecified MLE
selects the Poisson distribution closest to the true distribution
in Kullback--Leibler divergence. Matching its mean does not ensure
that its zero probability or variance is correct.

\subsection*{(d) Failure of the information identity}
For one observation, the Poisson score and Hessian are
\[
s_\lambda(y)=\frac{\partial\log f_\lambda(y)}{\partial\lambda}
=\frac y\lambda-1,\qquad
\frac{\partial^2\log f_\lambda(y)}{\partial\lambda^2}
=-\frac y{\lambda^2}.
\]
Under the true negative binomial distribution,
\[
J(\lambda)=\operatorname{Var}_{g_{r,p}}\{s_\lambda(y)\}
=\frac{r(1-p)}{p^2\lambda^2},\qquad
H(\lambda)=-E_{g_{r,p}}\left[
\frac{\partial^2\log f_\lambda(y)}{\partial\lambda^2}\right]
=\frac{r(1-p)}{p\lambda^2}.
\]
Thus $J(\lambda)=H(\lambda)/p>H(\lambda)$ for $0<p<1$:
the information equality fails.
At $\lambda^*=\mu$,
\[
J(\lambda^*)=\frac1{r(1-p)},\qquad
H(\lambda^*)=\frac p{r(1-p)}.
\]
The sandwich asymptotic variance of
$\sqrt n(\widehat\lambda-\lambda^*)$ is
\[
H(\lambda^*)^{-1}J(\lambda^*)H(\lambda^*)^{-1}
=\frac{r(1-p)}{p^2}=v.
\]
The inverse Hessian alone gives $H(\lambda^*)^{-1}=\mu=pv$,
understating the variance by the multiplicative factor $p$.
